% =====================================================================
%  03_instalacion.tex — Requisitos, configuración y despliegue
% =====================================================================
\chapter{Instalación y despliegue}
\label{cap:instalacion}

\section{Requisitos}

\subsection{Hardware}

\begin{itemize}
  \item NVIDIA Jetson Orin (probado en modo de potencia de 25 W).
  \item Cámara IP con salida RTSP, montada frente a la puerta del patio.
  \item Almacenamiento suficiente para 72 horas de video: alrededor de
        \textbf{65 GB} (unos 0.9 GB por hora).
  \item Red local entre la Jetson y la cámara.
\end{itemize}

\subsection{Software}

\begin{itemize}
  \item JetPack 7.2 con CUDA 13.2 y TensorRT 10.16 (es la base de la imagen).
  \item Docker y el plugin \texttt{docker compose}.
  \item Los modelos convertidos a TensorRT (\texttt{.engine}) en la carpeta
        \path{modelos/}. Sin ellos el sistema no arranca.
\end{itemize}

\section{Qué vive fuera de la imagen Docker}

El contenedor es desechable; lo que se conserva está en la Jetson:

\begin{longtable}{@{}L{4.5cm}L{10cm}@{}}
  \toprule
  \textbf{Ruta} & \textbf{Contenido} \\
  \midrule
  \endhead
  \path{.env} & Configuración y contraseñas. Único archivo que lee
    \path{config.py}. \\
  \path{datos/} & Base de datos SQLite (\texttt{detecciones.db}). \\
  \path{modelos/} & Motores TensorRT \texttt{.engine}. \\
  \path{~/Documents/Jetson8/} & Fotos, recortes, grabaciones y
    el video de prueba. \\
  \path{_archivo/} & Código viejo y respaldos. \textbf{No se usa.} \\
  \bottomrule
  \caption{Rutas persistentes en la Jetson.}
  \label{tab:rutas}
\end{longtable}

\begin{peligro}[title={Lo que nunca se sube a Git}]
\path{.env} (contraseñas), \path{whatsapp/sesion/} (equivale a la contraseña de
la cuenta de WhatsApp), \path{datos/} y los modelos \texttt{.engine}. Todos
están en el \path{.gitignore}.
\end{peligro}

\section{Configuración: el archivo \texttt{.env}}

Todas las variables están documentadas en \path{.env.example}, que sí se
versiona. El procedimiento es:

\begin{lstlisting}[language=bash, caption={Preparar la configuración},
                   label={lst:env}]
cd ~/Documents/ProyectoVisionCISA
cp .env.example .env
nano .env          # llenar RTSP_URL, modelos, Telegram y WhatsApp
\end{lstlisting}

Las variables \textbf{obligatorias} son las que \path{config.py} revisa al
arrancar; si falta una, el proceso se detiene con un mensaje claro:

\begin{center}
\begin{tabular}{@{}ll@{}}
  \toprule
  \textbf{Grupo} & \textbf{Variables obligatorias} \\
  \midrule
  Cámara & \texttt{RTSP\_URL} \\
  Modelos & \texttt{MODEL\_PATH}, \texttt{DIGITS\_MODEL\_PATH},
            \texttt{AUTOBUS\_MODEL\_PATH} \\
  Telegram & \texttt{TG\_TOKEN}, \texttt{TG\_CHAT\_ID}, \texttt{TG\_VALIDADOR\_IDS} \\
  Patio & \texttt{ID\_PUERTA}, \texttt{FRAMES\_DIR} \\
  Web & \texttt{FLASK\_HOST}, \texttt{FLASK\_PORT} \\
  \bottomrule
\end{tabular}
\end{center}

La referencia completa de variables, con valores por omisión y efecto, está en
la sección~\ref{sec:variables}.

\section{Arranque}

\begin{lstlisting}[language=bash, caption={Ciclo normal de operación},
                   label={lst:arranque}]
cd ~/Documents/ProyectoVisionCISA

docker compose up -d                 # arranca deteccion y grabacion
docker compose ps                    # cisa-vision y cisa-grabacion en Up
docker compose logs -f --tail 30     # mensajes en vivo (Ctrl+C para salir)
docker compose down                  # detiene todo; los datos NO se borran
\end{lstlisting}

Después de cambiar el \path{.env} hay que recrear el contenedor para que lo lea:

\begin{lstlisting}[language=bash]
docker compose up -d --force-recreate
\end{lstlisting}

Después de cambiar el código hay que reconstruir la imagen:

\begin{lstlisting}[language=bash]
docker compose build && docker compose up -d --force-recreate
\end{lstlisting}

\begin{aviso}[title={No arranca solo al prender la Jetson}]
En \path{docker-compose.yml} los servicios tienen \texttt{restart: "no"}. Si se
quiere que arranquen solos después de un corte de luz, cambia ese valor a
\texttt{unless-stopped} y vuelve a levantar los servicios.
\end{aviso}

\section{Verificación}

Con los servicios arriba, comprueba en este orden:

\begin{enumerate}
  \item \textbf{Contenedores}: \texttt{docker compose ps} muestra
        \texttt{cisa-vision} y \texttt{cisa-grabacion} en estado \texttt{Up}.
  \item \textbf{Salud del pipeline}: abre
        \url{http://<ip-jetson>:5001/health}; debe responder con el estado de la
        conexión, el FPS de YOLO y la cámara activa.
  \item \textbf{Tablero}: \url{http://<ip-jetson>:5001/}. Se ve el video y el
        contador de unidades dentro.
  \item \textbf{Video en vivo}: \url{http://<ip-jetson>:5001/livevideo}.
  \item \textbf{Base de datos}: \texttt{docker compose logs} no debe mostrar
        errores de SQLite ni de migración.
\end{enumerate}

\section{Modo prueba con video}

Para trabajar sin cámara, apunta \texttt{RTSP\_URL} a un archivo \texttt{.mp4}
y apaga los avisos:

\begin{lstlisting}[language=bash, caption={Cambiar entre cámara y video},
                   label={lst:fuente}]
python3 scripts/fuente.py video    # usa el .mp4 de prueba
python3 scripts/fuente.py camara   # regresa a la camara
python3 scripts/fuente.py          # sin argumento: dice cual esta activa
\end{lstlisting}

En el \path{.env}, \texttt{SIN\_TELEGRAM=1} evita mandar mensajes reales y
\texttt{VALIDAR\_CON\_GEMINI} controla la revisión automática.

\section{Actualizar el sistema}

\begin{enumerate}
  \item Respalda la base: copia \path{datos/detecciones.db}.
  \item Trae los cambios: \texttt{git pull}.
  \item Reconstruye: \texttt{docker compose build}.
  \item Levanta: \texttt{docker compose up -d --force-recreate}.
  \item Revisa los logs y \texttt{/health}.
\end{enumerate}

\begin{nota}[title={La base se migra sola}]
Al arrancar, \path{metricas/storage.py} aplica el esquema y, si encuentra una
base del formato anterior, la respalda y la migra. Aun así, un respaldo antes de
actualizar nunca sobra.
\end{nota}
